\title{DeepSeekMath: Pushing the Limits of Mathematical Reasoning in Open Language Models}

\begin{abstract}
The task of mathematical reasoning remains particularly demanding for language models, largely because of its inherently intricate and highly structured characteristics. Here, we present DeepSeekMath~7B, a model obtained by extending the pre-training of DeepSeek-Coder-Base-v1.5 7B with an additional 120B mathematics-focused tokens extracted from Common Crawl, supplemented by datasets containing natural language and code. DeepSeekMath~7B reaches a remarkable accuracy of 51.7\% on the advanced MATH benchmark, all without the use of supplementary toolkits or ensemble voting methods, achieving results comparable to Gemini-Ultra and GPT-4. Employing self-consistency evaluation across 64 samples, DeepSeekMath~7B attains 60.9\% on MATH. The strong mathematical reasoning exhibited by DeepSeekMath~stems from two central strategies: (1) the effective utilization of large-scale web data via a rigorously designed selection and filtering pipeline; (2) the adoption of Group Relative Policy Optimization (GRPO), an adaptation of Proximal Policy Optimization (PPO), which bolsters mathematical reasoning proficiency and concurrently streamlines the memory demands inherent to PPO.
\end{abstract}

\section{Introduction}

Large language models (LLMs) have dramatically transformed the landscape of mathematical reasoning within artificial intelligence, leading to noteworthy progress on quantitative reasoning benchmarks \citep{MATH} as well as those focused on geometry \citep{trinh2024solving}. Furthermore, these systems have demonstrated their utility in supporting humans tackling sophisticated mathematical challenges \citep{tao}. Nevertheless, state-of-the-art systems such as GPT-4 \citep{gpt4} and Gemini-Ultra \citep{gemini} remain closed-source, with publicly accessible open-source alternatives still falling substantially short in terms of performance.

This paper presents DeepSeekMath, a specialized language model engineered for mathematical reasoning that surpasses the performance of existing open-source models and comes close to matching GPT-4 on standard academic math assessments. Central to our approach is the construction of the DeepSeekMath Corpus, an extensive and high-fidelity dataset containing 120B tokens specifically focused on mathematical content. The dataset is curated from Common Crawl (CC) through the deployment of a classifier based on fastText \citep{joulin2016fasttext}. Initially, this classifier is trained with positive samples sourced from OpenWebMath \citep{openwebmath} alongside a varied collection of general web content used as negative samples. This model is subsequently employed to identify further mathematical texts from CC, which are then subjected to meticulous human annotation for increased precision. The resulting labeled data is used to refine and retrain the classifier, thereby enhancing its ability to identify high-quality mathematical data. Evaluations of this corpus demonstrate its robust quality; notably, our foundational model, DeepSeekMath-Base 7B, attains a score of 64.2\% on the GSM8K benchmark \citep{gsm8k} and 36.2\% on the advanced MATH dataset \citep{MATH}, outstripping the performance of Minerva 540B \citep{minerva}. Additionally, the multilingual nature of the DeepSeekMath Corpus yields notable improvements on Chinese mathematics benchmarks \citep{wei2023cmath,agieval}. We contend that our methodology for processing mathematical data lays the groundwork for future advancements in the field, with ample opportunity for further progress.

DeepSeekMath-Base is built upon the DeepSeek-Coder-Base-v1.5 7B model \citep{deepseek-coder}, as empirical evidence suggests that initializing from a model trained on code leads to superior outcomes compared to starting from a generic large language model. We also observe that mathematical training enhances not only the model’s mathematical competence but also its general reasoning, as demonstrated by improved results on MMLU \citep{mmlu} and BBH \citep{bbh} benchmarks.

Following pre-training, we fine-tune DeepSeekMath-Base with mathematical instruction datasets that include chain-of-thought reasoning \citep{cot}, program-of-thought reasoning \citep{pot,pal}, and reasoning with tool integration \citep{tora}. The resultant DeepSeekMath-Instruct 7B surpasses all other 7B models and performs on par with the leading 70B-scale open-source instruction-tuned models.

In addition, we propose Group Relative Policy Optimization (GRPO), a reinforcement learning method derived from Proximal Policy Optimization (PPO) \citep{schulman2017proximal}. GRPO dispenses with the use of a critic model, opting instead to estimate baselines using group scores, thus significantly lowering computational demands. Utilizing only a selected subset of English instruction tuning data, GRPO delivers marked improvements over the already robust DeepSeekMath-Instruct, yielding gains both for in-domain tasks (GSM8K: 82.9\% $\rightarrow$ 88.2\%, MATH: 46.8\% $\rightarrow$ 51.7\%) and out-of-domain mathematical benchmarks (e.g., CMATH: 84.6\% $\rightarrow$ 88.8\%) during the reinforcement learning stage. We also introduce a unified framework for understanding multiple approaches such as Rejection Sampling Fine-Tuning (RFT) \citep{yuan2023scaling}, Direct Preference Optimization (DPO) \citep{dpo}, PPO, and our proposed GRPO, showing that these techniques can be interpreted as direct or simplified forms of reinforcement learning. Through comprehensive experiments—including comparisons between online and offline training, evaluation of outcome versus process supervision, and analysis of single-turn versus iterative reinforcement learning—we explore the critical factors shaping this paradigm. Finally, we elucidate the mechanisms by which our RL approach enhances instruction-tuned model performance, and outline future research directions for refining reinforcement learning within this unified context.

\subsection{Contributions}
This work advances both large-scale mathematical pre-training and the study of reinforcement learning through comprehensive empirical analysis.

\textbf{Math Pre-Training at Scale}

\begin{itemize}[topsep=0pt]
    \item We demonstrate that Common Crawl’s publicly available dataset harbors substantial resources for mathematical modeling. Employing a rigorous data filtering and selection strategy, we curate the DeepSeekMath Corpus, a high-quality mathematical dataset comprising 120B tokens. This dataset is nearly 7-fold greater than that of Minerva’s math web set \citep{minerva} and roughly 9 times larger than the recently introduced OpenWebMath \citep{openwebmath}.

    \item The DeepSeekMath-Base 7B model we pre-train exhibits performance on par with Minerva 540B \citep{minerva}. Our results challenge the notion that parameter count is the primary determinant of mathematical reasoning in large language models, showing that compact models trained on targeted, high-quality datasets can perform competitively.

    \item We report experimental insights concerning math-specific training. Initial code pre-training, followed by mathematical data pre-training, substantially enhances models' mathematical problem-solving skills, regardless of tool usage. These results contribute partial evidence to the debated hypothesis: \textit{does code training boost reasoning skills?} Our evidence supports this, at least in the context of mathematical reasoning.

    \item While leveraging arXiv papers for pre-training is prevalent within the community, our empirical assessment indicates that their inclusion offers negligible gains on the mathematical benchmarks considered in this study.
\end{itemize}

\textbf{Exploration and Analysis of Reinforcement Learning}

\begin{itemize}[topsep=0pt]
    \item We present Group Relative Policy Optimization (GRPO), a novel and resource-efficient reinforcement learning approach. Distinct from conventional algorithms like Proximal Policy Optimization (PPO), GRPO eliminates the critic component, utilizing group-based baseline estimation. This adjustment substantially reduces computational demands during training.
    \item Through extensive experiments, we show that GRPO markedly improves the performance of our instruction-tuned model, DeepSeekMath-Instruct, leveraging only instruction-tuning datasets. Additionally, we note advancements in generalization to out-of-domain data throughout the reinforcement learning phase.
    \item We offer a comprehensive framework that contextualizes existing algorithms—including RFT, DPO, PPO, and GRPO—within a unified perspective. Our study encompasses rigorous comparisons across various settings, such as online versus offline learning, outcome-based versus process-based supervision, and single-step versus iterative reinforcement learning. This allows for an in-depth exploration of key factors underlying the unified paradigm.
    \item Utilizing our consolidated framework, we analyze the mechanisms that underlie the success of reinforcement learning approaches and identify several promising avenues to further enhance reinforcement learning techniques for LLMs.
\end{itemize}

\subsection{Summary of Evaluations and Metrics}

\begin{itemize}[topsep=0pt]
    \item \textbf{English and Chinese Mathematical Reasoning}: We thoroughly evaluate our models across a spectrum of English and Chinese mathematical benchmarks, ranging from elementary to university-level problems. The suite of English tasks includes GSM8K \citep{gsm8k}, MATH \citep{MATH}, SAT \citep{llemma}, OCW Courses \citep{minerva}, and MMLU-STEM \citep{mmlu}, while Chinese assessments comprise MGSM-zh \citep{mgsm}, CMATH \citep{wei2023cmath}, Gaokao-MathCloze \citep{agieval}, and Gaokao-MathQA \citep{agieval}. Our evaluation focuses on the models’ capabilities to independently construct textual solutions and to resolve problems programmatically via Python.

    On English mathematical challenges, DeepSeekMath-Base rivals the performance of the proprietary Minerva 540B model \citep{minerva}, outperforming all open-source foundational models—such as Mistral 7B \citep{mistral} and Llemma-34B \citep{llemma}—by notable margins, regardless of their math-specific pre-training status. In Chinese settings, DeepSeekMath-Base achieves superior results, a likely outcome of including high-quality multilingual data rather than adhering to prior methods \citep{minerva,llemma} which focused exclusively on English mathematical corpora. With the integration of instruction tuning and reinforcement learning, the enhanced DeepSeekMath-Instruct and DeepSeekMath-RL models demonstrate excellent results, surpassing 50\% accuracy on the highly competitive MATH dataset—setting a new milestone among open-source contributions.
\end{itemize}

\item \textbf{Formal Mathematics}: To assess DeepSeekMath-Base’s proficiency in transforming informal mathematical statements into formal proofs, we utilize the informal-to-formal theorem proving benchmark described in \citep{dsp_proof}, specifically employing the miniF2F dataset \citep{minif2f} in conjunction with the Isabelle proof assistant \citep{isabelle}. Our findings show that DeepSeekMath-Base achieves robust few-shot performance in the autoformalization setting.

\item \textbf{Natural Language Understanding, Reasoning, and Code}: For a holistic evaluation of the model’s abilities in general language comprehension, reasoning, and programming, we administer DeepSeekMath-Base on several established benchmarks: Massive Multitask Language Understanding (MMLU) \citep{mmlu}, encompassing 57 diverse multiple-choice disciplines; BIG-Bench Hard (BBH) \citep{bbh}, containing 23 complex tasks with a focus on multi-step reasoning; along with the HumanEval \citep{codex} and MBPP \citep{mbpp} benchmarks, both widely adopted for assessing code generation proficiency. Our results indicate that mathematical pre-training confers notable improvements in both linguistic and reasoning task performance.

\section{Math Pre-Training}

\subsection{Data Collection and Decontamination}
Here, we describe the methodology for assembling the DeepSeekMath Corpus from Common Crawl. Figure \ref{fig:data_collect} illustrates the iterative data gathering procedure, outlining how a large-scale mathematical dataset can be constructed systematically, beginning from an initial high-quality, domain-specific corpus. The pipeline is readily transferable to other fields, such as software development.

The process begins with OpenWebMath \citep{openwebmath}, a curated compilation of high-quality mathematical web material, selected as our seed corpus. Leveraging this initial set, we train a fastText classifier \citep{joulin2016fasttext} to identify and retrieve additional web documents resembling OpenWebMath. Specifically, the training uses 500,000 examples sampled from the seed corpus as positive instances and 500,000 randomly sampled Common Crawl web pages as negative instances. The fastText model is trained using the open-source implementation\footnote{\url{https://fasttext.cc}}, setting the embedding dimension to 256, the learning rate to 0.1, maximum n-gram length to 3, minimum frequency for word inclusion to 3, and performing 3 epochs. To further reduce redundancy in the dataset, URL-based deduplication and near-duplicate filtering are applied to the original Common Crawl, yielding approximately 40 billion unique HTML pages. We then apply the trained fastText classifier to identify mathematical documents from this deduplicated corpus. To ensure content quality, candidate documents are ranked by the classifier’s scores, with only the highest scoring samples retained. The optimal amount of data preserved is determined by conducting preliminary training runs on subsets of 40B, 80B, 120B, and 160B tokens; in the first iteration, we retain the top 40B tokens.

Following the initial round of data gathering, a significant portion of mathematical web pages remains uncaptured. This shortfall can primarily be attributed to the fastText model’s limited training data, which does not sufficiently represent the diversity present in mathematical websites. To address this, we expand the range of mathematical sources included in our seed dataset, enhancing the quality and breadth of training data for the fastText model. In particular, we start by segmenting the full Common Crawl dataset into non-overlapping domains—where a domain comprises all pages that share a single base URL. For each such domain, we compute the proportion of pages acquired during the first iteration. Those domains with over 10\% of their pages collected are tagged as math-related domains (for instance, \url{mathoverflow.net}). 

Next, within these math-related domains, we manually label the URLs that host mathematical material (e.g., \url{mathoverflow.net/questions}). Pages connected to these mathematical URLs but missing from our initial collection are appended to the seed corpus. This process expands our positive dataset and allows for re-training the fastText model with a more representative sample, which leads to improved coverage of mathematical web content in subsequent iterations. After conducting four rounds of collection, we amass a total of 35.5M mathematical pages, accumulating 120B tokens. Notably, in the fourth iteration, approximately 98\% of the pages overlap with those collected in the third round, prompting us to conclude the data gathering process.

To mitigate benchmark data leakage, we adopt the filtering protocols detailed in \cite{deepseek-coder}: we systematically exclude web pages that contain questions or solutions present in English mathematical evaluation sets such as GSM8K \citep{gsm8k} and MATH \citep{MATH}, as well as Chinese benchmarks such as CMATH \citep{wei2023cmath} and AGIEval \citep{agieval}. Specifically, any text segment exhibiting an exact 10-gram overlap with any substring from benchmark datasets is excluded from our math training set. For benchmark segments containing fewer than 10 but at least 3-grams, we similarly remove any matching text passages to prevent contamination.

\subsection{Evaluating the DeepSeekMath~Corpus Quality}

To assess the quality of the DeepSeekMath~Corpus, we perform pre-training studies that compare its efficacy against other prominent recently published mathematical corpora:

\begin{itemize}[topsep=0pt]
    \item \textbf{MathPile} \citep{mathpile}: an aggregated mathematical dataset (8.9B tokens) composed of sources such as textbooks, Wikipedia, ProofWiki, CommonCrawl, StackExchange, and arXiv, with the latter making up more than 85\% of the total corpus;
    \item \textbf{OpenWebMath} \citep{openwebmath}: a math-centric subset of CommonCrawl containing 13.6B tokens;
    \item \textbf{Proof-Pile-2} \citep{llemma}: a math-oriented dataset that integrates OpenWebMath, AlgebraicStack (10.3B tokens of mathematical code), and arXiv articles (28.0B tokens). For experiments on Proof-Pile-2, we adopt the 2:4:1 arXiv:Web:Code composition as recommended in \cite{llemma}.
\end{itemize}

\subsubsection{Training Setting}
\label{sec:quality-policy}
We perform mathematical domain adaptation on a base language model with 1.3B parameters, architecturally consistent with DeepSeek LLMs \citep{deepseek-llm} and referred to as DeepSeek-LLM 1.3B. For each distinct mathematical dataset, the model undergoes training for 150B tokens. All model training utilizes the streamlined and resource-efficient HAI-LLM \citep{haillm} training platform. Mirroring the training strategies of DeepSeek LLMs, we employ the AdamW optimizer \citep{adamW} with hyperparameters $\beta_1=0.9$, $\beta_2=0.95$, and $\mathrm{weight\_decay}=0.1$. A multi-step learning rate schedule is adopted: after 2,000 steps of linear warmup to the maximum rate (set at 5.3e-4), the rate drops to 31.6\% of its maximum after 80\% of total training steps, and finally to 10.0\% at the 90\% mark. During training, we use batches containing 4M tokens, each with a 4K sequence length.

\subsubsection{Evaluation Results}

\textbf{DeepSeekMath~Corpus delivers high quality, is multilingual, and constitutes the largest dataset in its class.}

\begin{itemize}[topsep=0pt]
    \item \textbf{High-quality}:
    We assess the model’s downstream abilities on eight different mathematical benchmarks, employing few-shot chain-of-thought prompting \cite{cot}. Table \ref{tab:corpora_comparison} indicates that the model fine-tuned on DeepSeekMath~Corpus consistently surpasses those trained on other corpora. As seen in Figure \ref{fig:corpora_comparisons}, even at 50B tokens—the equivalent of one full pass through Proof-Pile-2—the DeepSeekMath-trained model outperforms its counterpart, demonstrating that the overall data quality of DeepSeekMath~Corpus is superior.
    \item \textbf{Multilingual}:
    The DeepSeekMath~Corpus includes texts in various languages, with English and Chinese being the most prevalent. Table \ref{tab:corpora_comparison} reveals that fine-tuning on DeepSeekMath~Corpus substantially improves mathematical reasoning in both English and Chinese. Conversely, traditional corpora, which are heavily English-focused, exhibit limited transfer and may even negatively affect Chinese-language math tasks.
    \item \textbf{Large-scale}:
    DeepSeekMath~Corpus exceeds existing mathematical datasets by several factors in size. As displayed in Figure \ref{fig:corpora_comparisons}, DeepSeek-LLM 1.3B models trained on DeepSeekMath~Corpus not only learn faster but continue to make progress for a longer duration. By comparison, smaller baseline datasets are cycled through multiple times in training, after which the respective models' improvement stagnates.
\end{itemize}

\subsection{Training and Evaluating DeepSeekMath-Base 7B}

This section presents DeepSeekMath-Base 7B, a foundational model designed for robust mathematical reasoning tasks. The model's initialization draws from DeepSeek-Coder-Base-v1.5 7B \citep{deepseek-coder}, and it undergoes training over a corpus totaling 500B tokens. The data composition includes 56\% sourced from the DeepSeekMath Corpus, 4\% from AlgebraicStack, 10\% from arXiv, 20\% comprising code from Github, and the final 10\% consists of natural language samples from the Common Crawl in both English and Chinese. Training protocols largely mirror those described in Section \ref{sec:quality-policy}, with modifications: the peak learning rate is adjusted to 4.2e-4, and training utilizes a batch size of 10M tokens.

A thorough evaluation is performed to gauge the mathematical reasoning capabilities of DeepSeekMath-Base 7B. The analysis emphasizes its proficiency in independently generating mathematical solutions without external aid, problem-solving with the assistance of tools, and formal theorem proving. In addition to mathematics, we broaden the evaluation scope to profile the model’s effectiveness in natural language understanding, reasoning abilities, and programming tasks.

\paragraph{Mathematical Problem Solving with Step-by-Step Reasoning}

We assess the model's aptitude for addressing mathematical problems using a few-shot chain-of-thought prompting paradigm \citep{cot} across eight datasets in both English and Chinese. These benchmarks include a spectrum of quantitative reasoning datasets (such as GSM8K \citep{gsm8k}, MATH \citep{MATH}, and CMATH \citep{wei2023cmath}), as well as multiple-choice assessments (for example, MMLU-STEM \citep{mmlu} and Gaokao-MathQA \citep{agieval}), spanning a wide array of mathematical domains from foundational to advanced undergraduate material.

As illustrated in Table \ref{tab:base_math_cot}, DeepSeekMath-Base 7B achieves superior results across all eight evaluation benchmarks when compared to other open-source base models. This includes outperforming both the prominent general-purpose Mistral 7B model \citep{mistral} and Llemma 34B \citep{llemma}, which is a recent addition trained on the Proof-Pile-2 dataset \citep{llemma} with a focus on mathematics. In particular, for the challenging MATH competition benchmark, DeepSeekMath-Base records an improvement exceeding 10\% absolute over previous open-source base models, and also surpasses Minerva 540B \citep{minerva}—a closed-source base model built on the much larger PaLM architecture \citep{palm} that has undergone additional math-specific training.

\paragraph{Mathematical Problem Solving with Tool Use} We assess program-supported mathematical reasoning on GSM8K and MATH tasks by employing few-shot program-of-thought prompting strategies \citep{pot,pal}. In this setup, the models are instructed to tackle each problem via the construction of Python code, leveraging modules such as \textit{math} and \textit{sympy} for sophisticated computations. The final solution is taken as the output produced upon executing these programs. According to Table \ref{tab:base_math_pot_proof}, DeepSeekMath-Base 7B establishes a new state-of-the-art among open-source models, overtaking the earlier leader, Llemma 34B.

\paragraph{Formal Mathematics}

Automation of formal proofs serves to improve both the correctness and efficiency of mathematical argumentation, an area drawing increased scholarly focus. We conduct an evaluation of DeepSeekMath-Base 7B using the informal-to-formal proof generation task outlined in \citep{dsp_proof}. This involves generating a formal proof in response to an informal description, its formalized statement, and an accompanying informal proof. Experiments are carried out using the miniF2F \citep{minif2f} benchmark, which features Olympiad-level mathematical problems. For each problem, a formal proof in Isabelle is constructed via few-shot prompting. In line with the methodology from \cite{dsp_proof}, we prompt the model to output proof sketches, which are subsequently elaborated with the automated prover Sledgehammer \citep{sledgehammer} to complete missing steps. Table \ref{tab:base_math_pot_proof} evidences the strong capabilities of DeepSeekMath-Base 7B in the autoformalization of proofs.

\paragraph{Natural Language Understanding, Reasoning, and Code}

To examine natural language comprehension, reasoning ability, and coding proficiency, we employ MMLU \citep{mmlu} for understanding, BBH \citep{bbh} for reasoning, and HumanEval \citep{codex} alongside MBPP \citep{mbpp} for code-related tasks. As presented in Table \ref{tab:understanding_reasoning_code}, DeepSeekMath-Base 7B realizes marked improvements in both MMLU and BBH relative to its predecessor DeepSeek-Coder-Base-v1.5 \citep{deepseek-coder}, underscoring the benefit of mathematics-focused training on these dimensions. Moreover, thanks to the continued integration of code token training, DeepSeekMath-Base 7B preserves the code generation performance levels established by DeepSeek-Coder-Base-v1.5 for the coding benchmarks. In aggregate, DeepSeekMath-Base 7B significantly exceeds the performance of Mistral 7B \citep{mistral} across all three evaluated domains—reasoning and coding.

\section{Supervised Fine-Tuning}

\subsection{SFT Data Curation}

We developed a comprehensive dataset for mathematical instruction-tuning, incorporating both English and Chinese problem sets that span various domains within mathematics and encompass a wide spectrum of difficulty levels. Each problem is accompanied by solutions structured in chain-of-thought (CoT) \citep{cot}, program-of-thought (PoT) \citep{pot,pal}, or tool-integrated reasoning formats \citep{tora}. The final curated dataset consists of 776K training samples.

\begin{itemize}[topsep=0pt]
    \item \textbf{English mathematical datasets}: Our approach involves annotating GSM8K and MATH problem sets with tool-integrated solutions. We further include a selected portion of MathInstruct \citep{MathInstruct} and the training portion of Lila-OOD \citep{lila}, where solutions leverage either CoT or PoT reasoning. The English portion covers a broad array of mathematical subjects such as algebra, probability, number theory, calculus, and geometry.
    \item \textbf{Chinese mathematical datasets}: We assemble mathematical questions relevant to the Chinese K-12 curriculum, which span 76 distinct subfields like linear equations. Solutions are provided in both CoT and tool-integrated reasoning modalities.
\end{itemize}

\subsection{Training and Evaluating DeepSeekMath-Instruct 7B}

This section describes the training of DeepSeekMath-Instruct 7B, an instruction-tuned version based on the DeepSeekMath-Base model. During the training process, input examples are randomly concatenated until the maximum context window of 4K tokens is filled. Training proceeds for 500 steps using a batch size of 256 and a fixed learning rate of $5\times 10^{-5}$.

Model performance is assessed on mathematical benchmarks, both with and without the incorporation of external tools, across four quantitative reasoning datasets in English and Chinese. We compare our results with prominent state-of-the-art systems available at the time:

\begin{itemize}[topsep=0pt]
    \item \textbf{Closed-source models} comprise: (1) the GPT series, with GPT-4 \citep{gpt4} and GPT-4 Code Interpreter~\footnote{\url{https://openai.com/blog/chatgpt-plugins\#\#code-interpreter}} representing the most advanced, (2) Gemini Ultra and Pro \citep{gemini}, (3) Inflection-2 \citep{inflection-2}, (4) Grok-1~\footnote{\url{https://x.ai/model-card}}, along with contemporary Chinese-developed systems such as (5) Baichuan-3~\footnote{\url{https://www.baichuan-ai.com}}, and (6) the newest GLM-4~\footnote{\url{https://open.bigmodel.cn/dev/api\#glm-4}} from the GLM family \citep{glm}.
\end{itemize}

These models are developed for broad applicability, with most having undergone various alignment processes.

\item \textbf{Open-source models} consist of: general-purpose LLMs such as (1) DeepSeek-LLM-Chat 67B \citep{deepseek-llm}, (2) Qwen 72B \citep{qwen}, (3) SeaLLM-v2 7B \citep{seallm}, and (4) ChatGLM3 6B \citep{chatglm3}. Additionally, specialized models with mathematical enhancements include: (5) InternLM2-Math 20B~\footnote{\url{https://github.com/InternLM/InternLM-Math}}, which is built on InternLM2 and receives subsequent math-centric training followed by instruction tuning; (6) Math-Shepherd-Mistral 7B, where PPO training \citep{schulman2017proximal} is applied to Mistral 7B \citep{mistral} with a reward model based on process supervision; (7) the WizardMath family \citep{wizardmath}, which improves the mathematical reasoning of Mistral 7B and Llama-2 70B \citep{llama2} using evolve-instruct (an instruction tuning technique leveraging AI-generated instructions) and PPO training, with its training data primarily from GSM8K and MATH; (8) MetaMath 70B \citep{metamath}, a version of Llama-2 70B fine-tuned with an expanded version of GSM8K and MATH; (9) ToRA 34B \cite{tora}, a CodeLlama 34B model specifically trained for mathematical reasoning with integrated tool use; and (10) MAmmoTH 70B \citep{MathInstruct}, an instruction-tuned version of Llama-2 70B using the MathInstruct dataset.
\end{itemize}

Table \ref{tab:sft_rl_math} presents results for the evaluation setting in which tool use is not permitted. Here, DeepSeekMath-Instruct 7B displays strong capabilities in performing stepwise reasoning. In particular, on the challenging competition-level MATH benchmark, this model achieves scores at least 9\% higher (in absolute terms) than all other open-source models and most proprietary systems (such as Inflection-2 and Gemini Pro), including those with much larger parameter counts (e.g., Qwen 72B) or with explicit reinforcement learning geared towards mathematics (e.g., WizardMath-v1.1 7B). While DeepSeekMath-Instruct approaches the performance of the Chinese proprietary models GLM-4 and Baichuan-3 on MATH, it remains behind state-of-the-art systems such as GPT-4 and Gemini Ultra.

When problem solving allows the use of both natural language reasoning and code-based tool execution, DeepSeekMath-Instruct 7B reaches an accuracy near 60\% on MATH, outperforming all publicly available open-source models. On other evaluation benchmarks, it performs comparably to DeepSeek-LLM-Chat 67B, which previously led performance among open-source systems but is ten times larger.

\section{Reinforcement Learning}

\subsection{Group Relative Policy Optimization}
The use of reinforcement learning (RL) has proven effective for further enhancing the mathematical reasoning abilities of LLMs after the Supervised Fine-Tuning (SFT) phase \citep{wang2023math,wizardmath}. In this section, we present our RL method, Group Relative Policy Optimization (GRPO), designed to be both efficient and robust.

\subsubsection{From PPO to GRPO}
Proximal Policy Optimization (PPO) \citep{schulman2017proximal} is an actor-critic algorithm widely employed for RL fine-tuning of LLMs \citep{ouyang2022training}. The method tunes LLMs by optimizing the following surrogate loss function:

\begin{equation}
\footnotesize
    \mathcal{J}_{PPO}(\theta) = \mathbb{E}{[q \sim P(Q), o \sim \pi_{\theta_{old}}(O|q)]} \frac{1}{|o|} \sum_{t=1}^{|o|} \min \left[ \frac{\pi_\theta(o_{t} | q, o_{<t})}{\pi_{\theta_{old}}(o_{t} | q, o_{<t})} A_{t}, \text{clip} \left( \frac{\pi_\theta(o_{t} | q, o_{<t})}{\pi_{\theta_{old}}(o_{t} | q, o

Here, $\pi_{\theta}$ and $\pi_{\theta_{old}}$ denote the current and previous policy models, respectively, while $q$ and $o$ represent questions and outputs sampled from the dataset using the prior policy $\pi_{\theta_{old}}$. The hyper-parameter $\varepsilon$ is employed in PPO as part of a clipping strategy that enhances training stability. The advantage term $A_t$ is derived by employing Generalized Advantage Estimation (GAE) \citep{gae}, which utilizes reward signals $\{r_{\ge t}\}$ and an estimated value function $V_{\psi}$. As a result, PPO necessitates the simultaneous training of a value function along with the policy. To prevent excessive optimization toward the reward model, it is standard to incorporate a per-token KL penalty between the policy and a reference model into the reward for each token \citep{ouyang2022training}, formalized as follows:

\begin{equation}
    r_{t} = r_\varphi(q, o_{\le t}) - \beta \log\frac{\pi_{\theta}(o_{t}|q, o_{<t})}{\pi_{ref}(o_{t}|q, o_{<t})},
\label{eq:PPO-reward}
\end{equation}

where $r_\varphi$ stands for the reward model, $\pi_{ref}$ refers to the reference policy (commonly initialized as the SFT model), and $\beta$ is the scaling factor applied to the KL penalty.

Owing to the fact that PPO's value function often mirrors the policy model in scale, significant computational and memory resources are required. Additionally, in reinforcement learning, the value function acts as a baseline during advantage calculation to curtail variance. In the context of large language models, however, it is typical that the reward model assigns a score only to the final token, which introduces challenges for accurately training a token-level value function. To overcome this limitation, we introduce Group Relative Policy Optimization (GRPO), illustrated in Figure \ref{fig:grpo}, which eliminates the necessity for a separate value function. Instead, GRPO leverages the mean reward across multiple sampled outputs for each given question as its baseline. More concretely, for every question $q$, a set of outputs $\{o_1, o_2, \cdots, o_G\}$ is sampled from the prior policy $\pi_{\theta_{old}}$, and the policy model is updated to maximize the objective function:

\begin{equation}
\footnotesize
\begin{split}
    \mathcal{J}_{GRPO}(\theta) &= \mathbb{E}{[q \sim P(Q), \{o_i\}_{i=1}^G \sim \pi_{\theta_{old}}(O|q)]}  \\
    & \frac{1}{G}\sum_{i=1}^G\frac{1}{|o_i|} \sum_{t=1}^{|o_i|} \left\{ \min \left[ \frac{\pi_\theta(o_{i,t} | q, o_{i,<t})}{\pi_{\theta_{old}}(o_{i,t} | q, o_{i,<t})} \hat{A}_{i,t}, \text{clip} \left( \frac{\pi_\theta(o_{i,t} | q, o_{i,<t})}{\pi_{\theta_{old}}(o_{i,t} | q, o_{i,<t})}, 1 - \varepsilon, 1 + \varepsilon \right)  \hat{A}_{i,t} \right] - \beta \mathbb{D}_{KL}\left[\pi_{\theta} || \pi_{ref}\right]\right\} ,
\end{split}
\label{eq:GRPO-obj}
\end{equation}

Here, $\varepsilon$ and $\beta$ serve as hyper-parameters, while $\hat{A}_{i,t}$ denotes the advantage, which in this method is determined solely from the relative rewards within each group—a process elaborated on in subsequent sections. By utilizing relative rewards at the group level, GRPO is inherently compatible with reward models that are trained on pairwise or comparative judgments between outputs associated with the same query. Furthermore, rather than embedding the KL penalty within the reward computation as in PPO, GRPO directly introduces the KL divergence between the updated policy and the reference policy into the loss function, simplifying the process of determining $\hat{A}_{i,t}$. Notably, unlike the KL penalty applied in (\ref{eq:PPO-reward}), the KL divergence is here approximated using an unbiased estimator \citep{kl_approx}:

\begin{equation

\begin{algorithm}[t]
  \small
  \caption{Iterative Group Relative Policy Optimization}
  \textbf{Input} initial policy model $\pi_{\theta_{\text{init}}}$; reward models $r_\varphi$; task prompts $\mathcal{D}$; 
  hyperparameters $\varepsilon$, $\beta$, $\mu$
  \begin{algorithmic}[1]
    \State Initialize policy model $\pi_\theta \leftarrow \pi_{\theta_{\text{init}}}$
    \For{iteration = 1, \dots, I}
       \State Set reference model $\pi_{ref} \leftarrow \pi_{\theta}$
      \For{step = 1, \dots, M}
      \State Draw a batch $\mathcal{D}_b$ from $\mathcal{D}$
      \State Assign current policy model $\pi_{\theta_{old}} \leftarrow \pi_{\theta}$ 
      \State For each question $q \in \mathcal{D}_b$, sample $G$ outputs $\{o_i\}_{i=1}^G$ via $\pi_{\theta_{old}} (\cdot \mid q) $
      \State Compute rewards $\{r_i\}_{i=1}^{G}$ for generated outputs $o_i$ using $r_{\varphi}$ 
      \State Estimate group-relative advantages $\hat{A}_{i,t}$ for the $t$-th token in each $o_i$
      \For{GRPO iteration = 1, \dots, $\mu$}
        \State Update $\pi_{\theta}$ by maximizing the GRPO loss (Equation \ref{eq:GC-GRPO})
      \EndFor
    \EndFor \State Continuously update the reward model $r_\varphi$ through training with experience replay. 
    \EndFor 
  \end{algorithmic}
  \textbf{Output} $\pi_\theta$
  \label{alg:iter-grpo}
\end{algorithm}

\subsubsection{Outcome Supervision RL with GRPO} Formally, consider a prompt $q$ where $G$ outputs $\{o_1, o_2, \cdots, o_G\}$ are sampled from the previous policy $\pi_{\theta_{old}}$. Each output is evaluated by the reward model, producing $G$ associated scores $\mathbf{r}=\{r_1, r_2, \cdots, r_G\}$. To normalize the feedback, the mean reward of the group is subtracted and the result is divided by the group's standard deviation. Under outcome supervision, this normalized score is attributed to all tokens in a given output $o_i$, setting the tokenwise advantages as $\hat{A}_{i, t} = \widetilde{r}_i = \frac{r_i- {\rm mean}(\mathbf{r})}{{\rm std}(\mathbf{r})}$. The policy parameters are then adjusted to maximize the objective specified in equation (\ref{eq:GRPO-obj}).

\subsubsection{Process Supervision RL with GRPO} Since outcome supervision assigns rewards only at the sequence's conclusion, its signal may be insufficient for effectively shaping policies in challenging mathematical reasoning settings. Building on the approach of \cite{wang2023math}, we also investigate process supervision, in which feedback is provided after each intermediate reasoning stage. Given $q$ and its set of $G$ sampled outputs $\{o_1, o_2, \cdots, o_G\}$, a process-level reward model produces stepwise evaluations, leading to a reward structure $\mathbf{R} = \{ \{r_1^{index(1)},\cdots,r_1^{index(K_1)}\}, \cdots,  \{r_G^{index(1)},\cdots,r_G^{index(K_G)}\} \}$, where $index(j)$ denotes the endpoint token of step $j$ and $K_i$ is the number of reasoning steps in output $o_i$. These stepwise rewards are normalized using the mean and standard deviation across the group, such that $\widetilde{r}_i^{index(j)} = \frac{r_i^{index(j)} - {\rm mean(\mathbf{R})}}{{\rm std(\mathbf{R})}}$.
The token advantage $\hat{A}_{i, t}$ is defined as the sum of normalized rewards for all steps ending at or after token $t$, i.e., $\hat{A}_{i, t} = \sum_{index(j) \ge t} \widetilde{r}_i^{index(j)}$. The policy is optimized with respect to the objective in equation (\ref{eq

\subsubsection{Iterative RL with GRPO} As reinforcement learning proceeds, earlier versions of the reward model might not provide adequate supervision for the evolving policy model. To address this, we investigate iterative RL leveraging GRPO. As illustrated in Algorithm \ref{alg:iter-grpo}, this method involves periodically generating updated training datasets for the reward model, using samples drawn from the policy model’s current outputs. The reward model is updated in an ongoing manner utilizing a replay strategy that maintains 10\% of data from prior iterations. Subsequently, the latest policy model serves as the new reference model, and additional training is carried out on the policy model, now guided by the improved reward model.

\subsection{Training and Evaluating DeepSeekMath-RL}

Our reinforcement learning framework builds upon DeepSeekMath-Instruct 7B. The RL training dataset is curated from chain-of-thought-format problems related to GSM8K and MATH, extracted from the SFT data and amounting to roughly 144,000 instances. To assess the influence of RL on evaluation sets that do not have exposure during RL, all other SFT questions are omitted. The reward model’s training set is prepared as described in \citep{wang2023math}. Initial reward model training uses DeepSeekMath-Base 7B, employing a learning rate of 2e-5. For policy optimization with GRPO, the learning rate is fixed at 1e-6, with a KL regularization coefficient of 0.04. For every question, we produce $64$ completions. Maximum sequence length is limited to 1024 tokens, with a batch size of 1024 during training. The policy model receives only one gradient update following each sampling round. For evaluation, DeepSeekMath-RL 7B is assessed on the same benchmarks as DeepSeekMath-Instruct 7B. Within this setup, GSM8K and MATH, when presented with chain-of-thought reasoning, are considered in-domain, whereas all remaining benchmarks are treated as out-of-domain.

Table \ref{tab:sft_rl_math} presents comparative results for both open-source and proprietary models that utilize chain-of-thought and tool-integrated reasoning strategies on benchmarks in both English and Chinese. The analysis reveals the following insights: 1) When employing chain-of-thought reasoning, DeepSeekMath-RL 7B achieves 88.2\% accuracy on GSM8K and 51.7\% accuracy on MATH, outperforming all open-source models spanning the 7B to 70B parameter range and surpassing most closed-source systems as well. 2) Notably, DeepSeekMath-RL 7B is exclusively trained on chain-of-thought format instruction data from GSM8K and MATH, with initialization from DeepSeekMath-Instruct 7B. In spite of this focused training set, it consistently outperforms DeepSeekMath-Instruct 7B in every evaluation metric, thereby illustrating the advantages conferred by reinforcement learning.

\section{Discussion}
This section details key observations obtained from our pre-training and reinforcement learning experiments.

\subsection{Lessons Learnt in Pre-Training}

We begin by recounting insights drawn from the pre-training phase. Unless indicated otherwise, our experimental configurations align with those described in Section \ref{sec:quality-policy}. For clarity, in this discussion, the DeepSeekMath Corpus refers to the 89B-token collection assembled during the project's second data collection cycle.

\subsubsection{Code Training Benefits Mathematical Reasoning}

Although there is a widespread, yet often unsubstantiated, belief that code training enhances reasoning capability, our work aims to shed light on this assertion, particularly in mathematical contexts. Specifically, code-focused pre-training is found to bolster the capacity of models to perform mathematical reasoning, irrespective of tool assistance.

To investigate the impact of code-oriented training on mathematical reasoning, we conducted experiments comparing two distinct approaches: two-stage and one-stage training regimens.

\textbf{Two-Stage Training}

\begin{itemize}[topsep=0pt]
    \item \textbf{Code Training for 400B Tokens $\rightarrow$ Math Training for 150B Tokens}: We conduct sequential pretraining of DeepSeek-LLM 1.3B, first on 400B code tokens, then on an additional 150B math tokens;
    \item \textbf{General Training for 400B Tokens $\rightarrow$ Math Training for 150B Tokens}: As a baseline, we perform an analogous experiment using 400B general corpus tokens (sampled from DeepSeek-AI’s comprehensive general dataset) in the initial training phase in place of code, to directly compare the effect of code tokens versus general domain tokens in bolstering mathematical reasoning skills.
\end{itemize}

\textbf{One-Stage Training}

\begin{itemize}[topsep=0pt]
    \item \textbf{Math Training for 150B Tokens}: DeepSeek-LLM 1.3B is trained solely on 150B math tokens in a single phase;
    \item \textbf{Training on a mixture of 400B Code Tokens and 150B Math Tokens}: Since introducing a subsequent math training phase after prior code training impairs coding ability, we examine whether simultaneous exposure to a mixed dataset containing both 400B code and 150B math tokens in a unified training stage can simultaneously support mathematical reasoning and counteract catastrophic forgetting.
\end{itemize}

\paragraph{Results}
Tables \ref{tab:code-math} and \ref{tab:code-math-general-eval} provide empirical results comparing various training configurations on downstream tasks.

Introducing code pretraining augments program-aided mathematical problem solving, across both sequential and mixed training paradigms. Specifically, Table \ref{tab:code-math} indicates that code pretraining alone, under the sequential regime, leads to robust improvements in performance on Python-assisted tasks such as GSM8K and MATH, with further gains following subsequent math-focused finetuning. Notably, under the mixed regime, integrating code and math tokens during a single stage curbs the catastrophic forgetting observed in the sequential setting, enhancing both coding proficiency (as indicated in Table \ref{tab:code-math-general-eval}) and code-assisted mathematical reasoning (Table \ref{tab:code-math}).

Code pretraining further provides moderate boosts to mathematical reasoning tasks that do not leverage external tools. When code training is used as the initial stage, not only does it offer direct improvements, but it also primes the model to make subsequent math training more effective, yielding optimal results overall. However, the one-stage mixed approach appears to undermine non-tool mathematical reasoning, potentially because the limited parameter budget of DeepSeek-LLM 1.3B does not permit it to simultaneously master both programming and math data distributions.

\subsubsection{ArXiv Papers Appear to Contribute Little to Mathematical Reasoning Enhancement}

While ArXiv papers are frequently utilized within the pre-training datasets for mathematical language models \citep{minerva,gpt-f,llemma,mathpile}, their concrete effect on mathematical reasoning capabilities remains underexplored. Somewhat unexpectedly, our experimental findings indicate that ArXiv papers offer minimal benefit—if any—in fostering mathematical reasoning. We evaluate models of varying capacities, specifically DeepSeek-LLM 1.3B and DeepSeek-Coder-Base-v1.5 7B \citep{deepseek-coder}, each trained with different ArXiv-sourced datasets, which underwent distinct preprocessing methodologies:

\begin{itemize}[topsep=0pt]
    \item \textbf{MathPile} \citep{mathpile}: This dataset comprises 8.9B tokens, primarily cleaned and filtered using heuristic rules, and is made up of more than 85\% scientific ArXiv papers;
    \item \textbf{ArXiv-RedPajama} \citep{redpajama}: Consists of 28.0B tokens from the full set of ArXiv LaTeX documents, after exclusion of preambles, comments, macros, and bibliographies.
\end{itemize}

In our experimental protocol, DeepSeek-LLM 1.3B was trained for 150B tokens, and DeepSeek-Coder-Base-v1.5 7B for 40B tokens on each respective ArXiv dataset. Our results suggest that pre-training on ArXiv papers does not substantially boost mathematical reasoning performance. When the models are exposed exclusively to ArXiv-sourced data, we observe either stagnation or declines in performance across an array of mathematical tasks of varying levels of difficulty. Evaluated benchmarks include quantitative reasoning datasets such as GSM8K and MATH (Table \ref{tab:arxiv-cot}), multiple-choice STEM assessments like MMLU-STEM (Table \ref{tab:arxiv-cot}), and tasks in formal mathematics as represented by miniF2F (Table \ref{tab:arxiv_atp}).

It is important to acknowledge the boundaries of this conclusion. The current analysis does not address:

\begin{itemize}[topsep=0pt]
    \item The possible utility of ArXiv-sourced content for particular mathematical subtasks absent from our benchmarks, for instance, converting formal mathematical language to informal descriptions;
    \item The effects arising from mixing ArXiv data with alternative sources;
    \item Whether larger-scale models might realize benefits from ArXiv paper pre-training not visible at the current scales.
\end{itemize}

Consequently, further inquiry is warranted, and these questions are left open for subsequent research.

\subsection{Reflections on Reinforcement Learning}
\subsubsection{Toward a Unified Analytical Framework} Here, we introduce a generalized framework to comparatively assess diverse training strategies such as SFT, RFT, DPO, PPO, and GRPO, and proceed with experimental analysis to examine the components underlying this framework. Typically, the gradient of a training algorithm with respect to its parameters $\theta$ can be formulated as:

\begin{equation}
    \nabla_{\theta}\mathcal{J}_{\textcolor{red}{\mathcal{A}}}(\theta) = \mathbb{E}[\underbrace{(q,o) \sim \textcolor{red}{\mathcal{D}}}_{Data \ Source}]\left( \frac{1}{|o|} \sum_{t=1}^{|o|}  \underbrace{GC_{{\mathcal{A}}}(q, o, t, \textcolor{red}{\pi_{{rf}}})}_{Gradient \ Coefficient}  \nabla_{\theta}\log \pi_{\theta}(o_t | q, o_{<t})\right).
\label{eq:objective}
\end{equation}

This formulation contains three essential elements: 1) \textit{Data Source $\mathcal{D}$}, specifying the dataset utilized for training; 2) \textit{Reward Function $\pi_{{rf}}$}, which supplies the reward signal that guides learning; and 3) \textit{Algorithm $\mathcal{A}$}, which integrates the data and reward information to generate the gradient coefficient $GC$, thereby quantifying the degree of penalization or incentivization assigned to each sample. On the basis of this general framework, we review various representative strategies:

\begin{itemize}
    \item  \textbf{Supervised Fine-tuning (SFT)}: In SFT, the pretrained model undergoes further optimization on curated, human-annotated SFT datasets.
    \item \textbf{Rejection Sampling Fine-tuning (RFT)}: This approach further adapts the SFT model using a subset of responses sampled from the SFT model’s output, where outputs are filtered and selected according to the accuracy of the answers with respect to the SFT prompts.
    \item \textbf{Direct Preference Optimization (DPO)}: DPO takes the SFT model and enhances it via additional training on outputs generated by the SFT model, optimizing a loss function based on pairs of preferences.
    \item \textbf{Online Rejection Sampling Fine-tuning (Online RFT)}: In contrast to standard RFT, Online RFT uses the SFT model to initialize the policy, and subsequently improves it by fine-tuning with real-time outputs produced by the evolving policy model.
    \item \textbf{PPO/GRPO}: Here, the SFT model serves as the starting point for the policy, which is then updated and reinforced using trajectories sampled from the continuously updated policy itself.
\end{itemize}

A summary of the elements comprising these approaches is provided in Table \ref{tab:data_policy}. For an in-depth derivation, consult Appendix \ref{app:analysis-rl}.

\paragraph{Observations Regarding Data Source} We categorize data sources into online and offline sampling. Online sampling refers to situations where training samples are generated by the current, actively learning policy during its own exploration, whereas offline sampling refers to samples collected from the outputs of the initial SFT model. The RFT and DPO methods employ offline sampling, while Online RFT and GRPO employ online sampling.

Referencing Figure \ref{fig:rl-analysis}, it is evident that Online RFT yields considerably better outcomes than RFT across two evaluation benchmarks. During early training epochs, the performance of Online RFT and RFT is largely comparable; however, as training progresses, Online RFT obtains a clear performance lead. This finding underscores the advantages of online training. The rationale is that initially, the actor closely mimics the SFT model, and differences in sampled data are minimal. As training continues, however, samples obtained from the evolving actor diverge more substantially from the SFT, thereby increasing the benefits associated with sampling in real time.

\paragraph{Observations Regarding Gradient Coefficient} The reward signal is processed into a gradient coefficient for the purpose of updating model parameters. In our experimental design, we implement two kinds of reward functions: ‘Rule’ and ‘Model’. The ‘Rule’ approach scores responses according to predefined criteria based on answer correctness, while the ‘Model’ approach involves training a reward model that evaluates responses, with its own training data derived from rule-based judgments. A central distinction between GRPO and Online RFT, as shown in Equations \ref{eq:GC-RFT} and \ref{eq:GC-GRPO}, is that GRPO adjusts its gradient coefficient according to the reward value output by the reward model, thus providing variable encouragement or discouragement dependent on response quality. Online RFT, by comparison, does not have this adaptive property; it treats all correct responses with equivalent reinforcement and does not explicitly penalize incorrect responses.

As illustrated in Figure \ref{fig:rl-analysis}, GRPO outperforms online RFT, underscoring the effectiveness of modifying both positive and negative gradient coefficients. Moreover, the combination GRPO+PS achieves better results than GRPO+OS, suggesting that incorporating step-aware, fine-grained gradient coefficients is beneficial. Additionally, we examine the use of iterative RL; in our setup, we apply two iterations. Figure \ref{fig:iter-rl} reveals that this iterative process notably boosts performance, especially after the initial iteration.

\subsubsection{Why RL Works?} In this study, reinforcement learning is applied to a subset of the instruction tuning dataset, resulting in marked gains over the original instruction-tuned model. To better clarify the mechanisms behind these improvements, we assess both Pass@K and Maj@K accuracies of the Instruct and RL-enhanced models on two benchmark datasets. Figure \ref{fig:maj-pass} shows that RL boosts Maj@K scores, but does not yield improvements in Pass@K. These outcomes suggest that RL enhances the robustness of the model’s output distribution—specifically, \textbf{the observed improvement likely stems from elevating the correct response into the TopK set, rather than augmenting core reasoning capabilities.} Similarly, \citep{wang2023making} point out a \textbf{misalignment problem} in reasoning tasks within the SFT model, and demonstrate that SFT models can see their reasoning ability improved through various preference alignment approaches \citep{yuan2023rrhf,song2023preference,wang2023making}.

\subsubsection{How to Achieve More Effective RL?}
\label{sec:effec_rl}
We have shown that RL proves highly effective on mathematical reasoning tasks. Additionally, we introduce a unified perspective that frames various training methodologies under a general RL paradigm, wherein all approaches are interpreted as either direct applications or simplifications of RL. According to Equation \ref{eq:objective}, this paradigm consists of three primary components: Data Source, Algorithm, and Reward Function. We discuss several prospective avenues for further research related to each component.

\paragraph{Data Source} The data source constitutes the foundation for all training regimes. Within RL settings, it is understood as the collection of unlabeled questions accompanied by outputs sampled from the policy model. Our experiments exclusively utilize questions from the instruction tuning phase and outputs generated via basic nucleus sampling. We suspect that this restricted sampling might account for the fact that our RL method only lifts the Maj@K metric. Future work will investigate the application of our RL framework to out-of-distribution question prompts and incorporate \textbf{advanced sampling (decoding) strategies} such as tree-search-based methods \citep{yao2023tree}. Additionally, developments in \textbf{efficient inference techniques} \citep{xia-etal-2023-speculative,leviathan2023fast,kwon2023efficient,xia2024unlocking}, which impact the efficacy of model exploration, remain a critical area for future improvements.

\paragraph{Algorithms} Algorithms leverage both data and the feedback from the reward signal to adjust model parameters through the gradient coefficient. According to Equation \ref{eq:objective}, most current approaches inherently \textbf{TRUST} the output of the reward function to modulate the conditional probability of specific tokens. Nevertheless, it is impractical to guarantee the reliability of the reward signal at all times, particularly for highly sophisticated tasks. For instance, the PRM800K dataset \citep{lightman2023let}, which has been meticulously labeled by experienced annotators, still contains roughly 20\% inaccurate labels\footnote{\url{https://github.com/openai/prm800k/issues/12\#issuecomment-1728491852}}. Therefore, it becomes essential to investigate reinforcement learning algorithms designed to withstand the effects of noisy or unreliable reward signals. We anticipate that such \textbf{WEAK-TO-STRONG} \citep{burns2023weak} alignment techniques could fundamentally transform learning paradigms.

\paragraph{Reward Function} The reward function delivers the primary feedback signal necessary for training. In reinforcement learning, this function is often implemented as a neural reward model. We identify three crucial avenues for advancing reward models: 1) \textbf{Strengthening the generalization capacity of reward models.} Effective generalization is critical for managing inputs that lie outside the training distribution or are the product of sophisticated decoding strategies; failure in this regard might result in reinforcement learning that simply preserves the current distribution of LLMs rather than genuinely enhancing their core abilities; 2) \textbf{Capturing the uncertainty of reward models.} Representing uncertainty could serve as a connective element between less robust reward models and weak-to-strong learning frameworks; 3) \textbf{Developing high-quality process reward models efficiently} so as to yield granular feedback signals that can more accurately direct reasoning tasks \citep{lightman2023let,wang2023math}.

\section{Conclusion, Limitation, and Future Work}

In this work, we introduce DeepSeekMath, an open-source model that surpasses all comparable models in the competition-level MATH benchmark and achieves performance that closely approaches proprietary systems. The foundation of DeepSeekMath~is DeepSeek-Coder-v1.5 7B, which is continually trained over 500B tokens, notably including 120B mathematical tokens extracted from Common Crawl. Through detailed ablation experiments, we demonstrate that web pages constitute a valuable resource for sourcing high-quality mathematical content, whereas data from arXiv contributes less effectively than anticipated. We also propose Group Relative Policy Optimization (GRPO), a modified form of Proximal Policy Optimization (PPO), that significantly enhances mathematical reasoning skills while reducing memory requirements. Our empirical evaluation indicates that GRPO remains beneficial even as DeepSeekMath-Instruct 7B achieves high benchmark scores. Furthermore, we outline a coherent framework for unifying various methodologies and discuss several promising avenues for developing more efficient reinforcement learning approaches.

Despite these advances, DeepSeekMath~exhibits some weaknesses relative to closed models, particularly in areas such as geometry and theorem proving. During initial testing, for example, the model struggled with questions involving geometric figures like triangles and ellipses, which suggests the presence of selection biases within both the pre-training and fine-tuning datasets. Moreover, due to limitations in model size, DeepSeekMath~falls short of GPT-4 in terms of few-shot learning; while GPT-4 benefits substantially from few-shot examples, DeepSeekMath~displays similar outcomes in both zero-shot and few-shot scenarios. Looking ahead, we aim to refine our data selection process to curate higher-quality pre-training corpora and will further investigate the approaches discussed in Section \ref{sec:effec_rl} for enhancing the reinforcement learning of large language models (LLMs).

\end{document}